\vspace{0.5em}\noindent\textbf{XÂY DỰNG HỆ THỐNG CHAT REALTIME CÓ KHẢ NĂNG MỞ RỘNG VỚI
SOCKET.IO, REDIS VÀ
DYNAMODB}\par

Chat realtime hiện nay xuất hiện trong rất nhiều hệ thống như thương mại
điện tử, tuyển dụng, chăm sóc khách hàng, mạng xã hội và các nền tảng
cộng tác.

Ở mức cơ bản, chúng ta chỉ cần một server để nhận và gửi tin nhắn. Tuy
nhiên, khi số lượng người dùng tăng lên, hệ thống bắt đầu gặp những câu
hỏi khó hơn:

\begin{itemize}
\tightlist
\item
  Làm thế nào để tin nhắn xuất hiện ngay lập tức?
\item
  Điều gì xảy ra khi người gửi và người nhận kết nối vào hai server khác
  nhau?
\item
  Tin nhắn được lưu ở đâu để không bị mất khi server khởi động lại?
\item
  Làm thế nào để tăng số lượng server mà không làm gián đoạn người dùng?
\end{itemize}

Một kiến trúc phổ biến để giải quyết những vấn đề này là kết hợp
Socket.IO, Redis và Amazon DynamoDB.

\vspace{0.5em}\noindent\textbf{1. Vì sao REST API chưa đủ cho chat
realtime?}\par

REST API hoạt động theo mô hình request-response:

\begin{Shaded}
\begin{Highlighting}[]
\NormalTok{Client gửi request}
\NormalTok{v}
\NormalTok{Server xử lý}
\NormalTok{v}
\NormalTok{Server trả response}
\end{Highlighting}
\end{Shaded}

Nếu dùng REST API cho chat, frontend thường phải liên tục gọi API để
kiểm tra tin nhắn mới:

\begin{Shaded}
\begin{Highlighting}[]
\NormalTok{GET /messages}
\NormalTok{Chờ vài giây}
\NormalTok{GET /messages}
\NormalTok{Chờ vài giây}
\NormalTok{GET /messages}
\end{Highlighting}
\end{Shaded}

Cách này gọi là polling.

Polling tương đối đơn giản nhưng có một số hạn chế:

\begin{itemize}
\tightlist
\item
  Tin nhắn có độ trễ phụ thuộc vào chu kỳ gọi API.
\item
  Client gửi request ngay cả khi không có tin nhắn mới.
\item
  Backend phải xử lý nhiều request không cần thiết.
\item
  Khi số lượng người dùng tăng, tải hệ thống cũng tăng nhanh.
\end{itemize}

Đối với hệ thống cần phản hồi gần như tức thời, WebSocket hoặc một thư
viện hỗ trợ realtime như Socket.IO thường phù hợp hơn.

\vspace{0.5em}\noindent\textbf{2. Socket.IO được dùng để giao tiếp
realtime}\par

Socket.IO cho phép duy trì kết nối hai chiều giữa client và server.

Sau khi kết nối được thiết lập, client và server có thể chủ động gửi sự
kiện cho nhau mà không cần tạo một request HTTP mới cho mỗi lần trao
đổi.

Luồng gửi tin nhắn có thể hiểu đơn giản như sau:

\begin{Shaded}
\begin{Highlighting}[]
\NormalTok{Người dùng A gửi tin nhắn}
\NormalTok{v}
\NormalTok{Socket.IO server nhận sự kiện}
\NormalTok{v}
\NormalTok{Server xử lý tin nhắn}
\NormalTok{v}
\NormalTok{Server gửi sự kiện mới}
\NormalTok{v}
\NormalTok{Người dùng B nhận tin nhắn}
\end{Highlighting}
\end{Shaded}

Một số sự kiện thường gặp:

\begin{itemize}
\tightlist
\item
  \texttt{conversation:join}
\item
  \texttt{message:send}
\item
  \texttt{message:new}
\item
  \texttt{user:typing}
\item
  \texttt{message:read}
\end{itemize}

Socket.IO còn hỗ trợ:

\begin{itemize}
\tightlist
\item
  Tự động kết nối lại khi mạng bị gián đoạn.
\item
  Chia người dùng vào các room.
\item
  Gửi sự kiện đến một người hoặc một nhóm người.
\item
  Fallback khi WebSocket không khả dụng.
\end{itemize}

Tuy nhiên, Socket.IO chỉ giải quyết tốt bài toán realtime khi hệ thống
có một server. Khi ứng dụng chạy nhiều server, chúng ta cần thêm một cơ
chế đồng bộ.

\vspace{0.5em}\noindent\textbf{3. Vấn đề xuất hiện khi chat service chạy nhiều
instance}\par

Giả sử hệ thống chỉ có một chat server:

\begin{Shaded}
\begin{Highlighting}[]
\NormalTok{Người dùng A {-}+}
\NormalTok{              +{-}{-} Chat server}
\NormalTok{Người dùng B {-}+}
\end{Highlighting}
\end{Shaded}

Cả hai người dùng đều kết nối vào cùng một server. Khi người dùng A gửi
tin nhắn, server có thể gửi trực tiếp đến kết nối của người dùng B.

Khi số lượng người dùng tăng, hệ thống có thể phải chạy nhiều server:

\begin{Shaded}
\begin{Highlighting}[]
\NormalTok{Chat server 1}
\NormalTok{Chat server 2}
\NormalTok{Chat server 3}
\end{Highlighting}
\end{Shaded}

Load balancer có thể phân phối người dùng như sau:

\begin{Shaded}
\begin{Highlighting}[]
\NormalTok{Người dùng A $\textbackslash{}rightarrow$ Chat server 1}
\NormalTok{Người dùng B $\textbackslash{}rightarrow$ Chat server 2}
\end{Highlighting}
\end{Shaded}

Lúc này, Chat server 1 không trực tiếp quản lý kết nối của người dùng B.

Nếu không có cơ chế liên lạc giữa các server, tin nhắn của người dùng A
có thể được lưu thành công nhưng người dùng B không nhận được thông báo
realtime.

Đây là vấn đề mà Redis có thể giải quyết.

\vspace{0.5em}\noindent\textbf{4. Redis giúp đồng bộ sự kiện giữa các chat
server}\par

Redis là một hệ thống lưu trữ dữ liệu trong bộ nhớ có tốc độ cao. Trong
kiến trúc chat realtime, Redis thường được sử dụng cho cơ chế
publish/subscribe.

Luồng hoạt động:

\begin{Shaded}
\begin{Highlighting}[]
\NormalTok{Người dùng A}
\NormalTok{v}
\NormalTok{Chat server 1}
\NormalTok{v Publish}
\NormalTok{Redis}
\NormalTok{v Subscribe}
\NormalTok{Chat server 2}
\NormalTok{v}
\NormalTok{Người dùng B}
\end{Highlighting}
\end{Shaded}

Khi Chat server 1 nhận được tin nhắn:

\begin{itemize}
\tightlist
\item
  Server xử lý và lưu tin nhắn.
\item
  Server phát sự kiện lên Redis.
\item
  Các chat server khác nhận được sự kiện.
\item
  Server đang giữ kết nối của người nhận gửi tin nhắn xuống client.
\end{itemize}

Có thể hiểu Redis giống như một ``kênh phát thanh'' dùng chung giữa các
chat server.

Redis phù hợp với nhiệm vụ này vì:

\begin{itemize}
\tightlist
\item
  Tốc độ xử lý cao.
\item
  Hỗ trợ publish/subscribe.
\item
  Phù hợp cho dữ liệu và sự kiện ngắn hạn.
\item
  Giúp các server hoạt động độc lập nhưng vẫn trao đổi được với nhau.
\end{itemize}

Redis trong trường hợp này không nhất thiết là nơi lưu lịch sử chat lâu
dài. Nhiệm vụ đó nên được giao cho một database bền vững hơn.

\vspace{0.5em}\noindent\textbf{5. DynamoDB được dùng để lưu trữ dữ liệu
chat}\par

Amazon DynamoDB là dịch vụ NoSQL database được quản lý bởi AWS.

Dữ liệu chat thường có một số đặc điểm:

\begin{itemize}
\tightlist
\item
  Số lượng tin nhắn tăng liên tục.
\item
  Tốc độ ghi có thể cao.
\item
  Tin nhắn thường được truy vấn theo cuộc hội thoại.
\item
  Ít cần các phép join phức tạp.
\item
  Nhu cầu sử dụng có thể tăng đột biến.
\end{itemize}

DynamoDB phù hợp với mô hình này vì có khả năng mở rộng tốt và không yêu
cầu người dùng tự quản lý server database.

Một thiết kế đơn giản cho bảng tin nhắn có thể sử dụng:

\begin{itemize}
\tightlist
\item
  Partition key: \texttt{conversation\_\allowbreak{}id}
\item
  Sort key: \texttt{timestamp} hoặc \texttt{message\_\allowbreak{}id}
\end{itemize}

Trong đó:

\begin{itemize}
\tightlist
\item
  \texttt{conversation\_\allowbreak{}id} dùng để nhóm các tin nhắn của cùng một cuộc
  hội thoại.
\item
  \texttt{timestamp} dùng để sắp xếp tin nhắn theo thời gian.
\end{itemize}

Khi cần lấy lịch sử chat, hệ thống truy vấn các bản ghi có cùng
\texttt{conversation\_\allowbreak{}id}.

Vai trò của Redis và DynamoDB có thể phân biệt như sau:

\begin{Shaded}
\begin{Highlighting}[]
\NormalTok{Redis}
\NormalTok{$\textbackslash{}rightarrow$ Truyền sự kiện realtime giữa các server}

\NormalTok{DynamoDB}
\NormalTok{$\textbackslash{}rightarrow$ Lưu trữ lịch sử tin nhắn}
\end{Highlighting}
\end{Shaded}

\vspace{0.5em}\noindent\textbf{6. Luồng gửi tin nhắn hoàn
chỉnh}\par

Một tin nhắn có thể đi qua hệ thống theo các bước sau:

\begin{enumerate}
\def\labelenumi{\arabic{enumi}.}
\tightlist
\item
  Người dùng nhập nội dung tin nhắn.
\item
  Client gửi sự kiện qua Socket.IO.
\item
  Chat server xác thực người dùng.
\item
  Server kiểm tra quyền tham gia cuộc hội thoại.
\item
  Tin nhắn được lưu vào DynamoDB.
\item
  Server publish sự kiện qua Redis.
\item
  Chat server đang giữ kết nối của người nhận nhận sự kiện.
\item
  Server gửi tin nhắn mới đến client của người nhận.
\end{enumerate}

Sơ đồ tổng quát:

\begin{Shaded}
\begin{Highlighting}[]
\NormalTok{Client A}
\NormalTok{|}
\NormalTok{| Socket.IO}
\NormalTok{v}
\NormalTok{Chat server 1}
\NormalTok{|}
\NormalTok{+{-}{-}{-}{-} Lưu dữ liệu {-}{-}{-}{-}{-}\textgreater{} DynamoDB}
\NormalTok{|}
\NormalTok{+{-}{-}{-}{-} Publish event {-}{-}{-}\textgreater{} Redis}
\NormalTok{                         |}
\NormalTok{                         v}
\NormalTok{                  Chat server 2}
\NormalTok{                         |}
\NormalTok{                         v}
\NormalTok{                      Client B}
\end{Highlighting}
\end{Shaded}

Kiến trúc này giúp hệ thống đạt được hai mục tiêu:

\begin{itemize}
\tightlist
\item
  Phản hồi realtime nhờ Socket.IO và Redis.
\item
  Không mất lịch sử tin nhắn nhờ DynamoDB.
\end{itemize}

\vspace{0.5em}\noindent\textbf{7. Kubernetes hỗ trợ scale chat service như thế
nào?}\par

Khi ứng dụng được đóng gói bằng container, Kubernetes có thể chạy nhiều
bản sao của chat service.

\begin{Shaded}
\begin{Highlighting}[]
\NormalTok{Chat service}
\NormalTok{+{-}{-} Pod 1}
\NormalTok{+{-}{-} Pod 2}
\NormalTok{+{-}{-} Pod 3}
\NormalTok{+{-}{-} Pod N}
\end{Highlighting}
\end{Shaded}

Kubernetes Service hoặc load balancer phân phối kết nối đến các pod.

Kubernetes còn hỗ trợ:

\begin{itemize}
\tightlist
\item
  Tự động khởi động lại container khi xảy ra lỗi.
\item
  Rolling update khi triển khai phiên bản mới.
\item
  Tăng hoặc giảm số lượng pod.
\item
  Health check thông qua liveness và readiness probe.
\item
  Phân phối traffic giữa nhiều pod.
\item
  Horizontal Pod Autoscaler dựa trên CPU, memory hoặc custom metrics.
\end{itemize}

Tuy nhiên, Kubernetes chỉ giúp chạy và quản lý nhiều instance. Việc đồng
bộ realtime giữa các instance vẫn cần Redis hoặc một message broker phù
hợp.

\vspace{0.5em}\noindent\textbf{8. Sticky session và Redis giải quyết hai vấn đề khác
nhau}\par

Khi Socket.IO sử dụng HTTP long polling, một client có thể tạo nhiều
request liên tiếp trong cùng một phiên.

Nếu các request bị chuyển đến nhiều server khác nhau, phiên kết nối có
thể gặp lỗi.

Sticky session giúp các request của cùng một client tiếp tục được chuyển
về cùng một server:

\begin{Shaded}
\begin{Highlighting}[]
\NormalTok{Client A $\textbackslash{}rightarrow$ Chat server 1}
\NormalTok{Client A $\textbackslash{}rightarrow$ Chat server 1}
\NormalTok{Client A $\textbackslash{}rightarrow$ Chat server 1}
\end{Highlighting}
\end{Shaded}

Tuy nhiên, sticky session không thay thế Redis.

Có thể phân biệt:

\begin{itemize}
\tightlist
\item
  Sticky session giữ kết nối của một client ổn định trên một server.
\item
  Redis truyền sự kiện giữa các server khác nhau.
\end{itemize}

Ngay cả khi đã bật sticky session, người gửi và người nhận vẫn có thể
nằm trên hai server khác nhau.

\vspace{0.5em}\noindent\textbf{9. Triển khai trên AWS}\par

Khi triển khai trên AWS, kiến trúc có thể sử dụng:

\begin{Shaded}
\begin{Highlighting}[]
\NormalTok{Application Load Balancer}
\NormalTok{v}
\NormalTok{Amazon EKS hoặc container service}
\NormalTok{v}
\NormalTok{Nhiều chat service instance}
\NormalTok{+{-}{-} Amazon DynamoDB}
\NormalTok{+{-}{-} Amazon ElastiCache/Valkey}
\end{Highlighting}
\end{Shaded}

Vai trò của các dịch vụ:

\begin{itemize}
\tightlist
\item
  Amazon EKS: chạy và quản lý các container chat service.
\item
  Application Load Balancer: nhận và phân phối kết nối từ người dùng.
\item
  Amazon DynamoDB: lưu dữ liệu chat.
\item
  Amazon ElastiCache hoặc Valkey: cung cấp Redis-compatible service.
\item
  Amazon CloudWatch: lưu log, metrics và cảnh báo.
\end{itemize}

Với project nhỏ hoặc môi trường thử nghiệm, có thể chạy Redis và
DynamoDB Local bằng Docker trước khi sử dụng dịch vụ AWS thật.

\vspace{0.5em}\noindent\textbf{10. Những lỗi thường gặp khi xây dựng chat
realtime}\par

\vspace{0.5em}\noindent\textbf{Chỉ lưu tin nhắn trong bộ nhớ của
server}\par

Nếu server restart, toàn bộ dữ liệu có thể bị mất.

\begin{Shaded}
\begin{Highlighting}[]
\NormalTok{Server restart $\textbackslash{}rightarrow$ Message history mất}
\end{Highlighting}
\end{Shaded}

Tin nhắn cần được lưu vào một database bền vững.

\vspace{0.5em}\noindent\textbf{Chạy nhiều server nhưng không có Redis
adapter}\par

Mỗi server chỉ biết các client đang kết nối trực tiếp với nó. Người dùng
ở server khác có thể không nhận được sự kiện.

\vspace{0.5em}\noindent\textbf{Sử dụng Redis làm nơi lưu trữ duy
nhất}\par

Redis phù hợp cho cache và truyền sự kiện, nhưng không phải lúc nào cũng
phù hợp để trở thành nguồn lưu trữ lịch sử chat duy nhất.

\vspace{0.5em}\noindent\textbf{Không xác thực Socket.IO
connection}\par

WebSocket hoặc Socket.IO connection vẫn cần được xác thực bằng JWT,
session hoặc một cơ chế phù hợp.

\vspace{0.5em}\noindent\textbf{Không kiểm tra quyền truy cập
conversation}\par

Một người dùng đã đăng nhập không có nghĩa là họ được phép tham gia mọi
phòng chat.

\vspace{0.5em}\noindent\textbf{Không xử lý tin nhắn gửi
trùng}\par

Khi mạng không ổn định, client có thể gửi lại cùng một message. Có thể
sử dụng client-generated message ID hoặc idempotency key để hạn chế bản
ghi trùng.

\vspace{0.5em}\noindent\textbf{Không giới hạn kích thước và tần suất gửi tin
nhắn}\par

Hệ thống nên có:

\begin{itemize}
\tightlist
\item
  Giới hạn độ dài message.
\item
  Rate limiting.
\item
  Chống spam.
\item
  Kiểm tra loại file nếu hỗ trợ attachment.
\end{itemize}

\vspace{0.5em}\noindent\textbf{11. Một số cách tối ưu chi
phí}\par

Hệ thống chat realtime trên AWS có thể phát sinh chi phí từ compute,
load balancer, database, Redis và log.

Một số cách tối ưu:

\begin{itemize}
\tightlist
\item
  Sử dụng DynamoDB On-Demand khi traffic còn thấp hoặc khó dự đoán.
\item
  Chọn ElastiCache node phù hợp thay vì sử dụng cấu hình quá lớn.
\item
  Dùng Auto Scaling để tránh chạy quá nhiều instance khi traffic thấp.
\item
  Đặt thời gian lưu CloudWatch Logs thay vì giữ log vô thời hạn.
\item
  Sử dụng DynamoDB TTL cho dữ liệu tạm thời nếu nghiệp vụ cho phép.
\item
  Kiểm tra và xóa Load Balancer, ElastiCache cluster hoặc môi trường thử
  nghiệm sau khi hoàn thành lab.
\item
  Chạy hệ thống bằng Docker và DynamoDB Local trong giai đoạn phát
  triển.
\end{itemize}

Việc sử dụng managed services giúp giảm công sức vận hành, nhưng vẫn cần
theo dõi tài nguyên để tránh phát sinh chi phí ngoài dự kiến.

\vspace{0.5em}\noindent\textbf{12. Checklist tham khảo}\par

\begin{itemize}
\tightlist
\item
  Sử dụng Socket.IO hoặc WebSocket cho giao tiếp realtime.
\item
  Dùng Redis adapter khi chạy nhiều chat server.
\item
  Lưu lịch sử tin nhắn trong database bền vững.
\item
  Thiết kế partition key DynamoDB phù hợp.
\item
  Không lưu state quan trọng chỉ trong memory của server.
\item
  Cấu hình sticky session nếu sử dụng long polling.
\item
  Xác thực kết nối realtime.
\item
  Kiểm tra quyền tham gia cuộc hội thoại.
\item
  Thêm rate limiting và giới hạn kích thước tin nhắn.
\item
  Cấu hình liveness và readiness probe.
\item
  Theo dõi số connection, message rate, latency và error rate.
\item
  Thiết lập Auto Scaling và resource limit phù hợp.
\item
  Kiểm tra chi phí các tài nguyên AWS sau khi thử nghiệm.
\end{itemize}

Bài học quan trọng nhất mình rút ra là: chat realtime không chỉ là thiết
lập một kết nối WebSocket.

Khi hệ thống bắt đầu mở rộng, chúng ta cần tách rõ ba nhiệm vụ:

\begin{Shaded}
\begin{Highlighting}[]
\NormalTok{Socket.IO}
\NormalTok{$\textbackslash{}rightarrow$ Giao tiếp realtime giữa client và server}

\NormalTok{Redis}
\NormalTok{$\textbackslash{}rightarrow$ Đồng bộ sự kiện giữa nhiều server}

\NormalTok{DynamoDB}
\NormalTok{$\textbackslash{}rightarrow$ Lưu trữ dữ liệu chat lâu dài}
\end{Highlighting}
\end{Shaded}

Việc lựa chọn đúng công cụ cho từng nhiệm vụ giúp hệ thống vừa phản hồi
nhanh, vừa có khả năng mở rộng khi số lượng người dùng tăng.

\texttt{\#AWS} \texttt{\#SocketIO} \texttt{\#Redis} \texttt{\#DynamoDB}
\texttt{\#AmazonEKS} \texttt{\#RealtimeChat} \texttt{\#NodeJS}
\texttt{\#Kubernetes} \texttt{\#CloudComputing} \texttt{\#AWSStudyGroup}
